\documentclass[10pt, a4paper]{article}

% --- Packages ---
\usepackage[utf8]{inputenc}
\usepackage[T1]{fontenc}
\usepackage[margin=0.60in, top=0.72in, bottom=0.72in]{geometry}
\usepackage{xcolor}
\usepackage{titlesec}
\usepackage{booktabs}
\usepackage{tabularx}
\usepackage{longtable}
\usepackage{enumitem}
\usepackage{amsmath, amssymb}
\usepackage{fancyhdr}
\usepackage{microtype}
\usepackage{multicol}
\usepackage{mdframed}
\usepackage[hidelinks]{hyperref}

% --- Monochromatic Color Palette ---
\definecolor{monodark}{HTML}{111111}
\definecolor{monomedium}{HTML}{3A3A3A}
\definecolor{monolight}{HTML}{666666}
\definecolor{monoline}{HTML}{CCCCCC}
\definecolor{monoshade}{HTML}{F0F0F0}
\definecolor{monobox}{HTML}{E8E8E8}

% --- Typography & Spacing ---
\setlist[itemize]{noitemsep, topsep=2pt, leftmargin=14pt}
\setlist[enumerate]{noitemsep, topsep=2pt, leftmargin=14pt}
\emergencystretch=4em

% --- Code Box Style ---
\newmdenv[
  backgroundcolor=monoshade,
  linecolor=monoline,
  linewidth=0.5pt,
  topline=true,
  bottomline=true,
  leftline=true,
  rightline=true,
  innerleftmargin=6pt,
  innerrightmargin=6pt,
  innertopmargin=4pt,
  innerbottommargin=4pt,
  skipabove=4pt,
  skipbelow=4pt,
]{codebox}

% --- Section Headers ---
\titleformat{\section}
  {\color{monodark}\large\bfseries}
  {\thesection.}{0.5em}{}
  [\color{monoline}\titlerule]

\titleformat{\subsection}
  {\color{monomedium}\normalsize\bfseries}
  {\thesubsection}{0.5em}{}

\titleformat{\subsubsection}
  {\color{monolight}\small\bfseries\itshape}
  {}{0em}{}

\titlespacing*{\section}{0pt}{10pt}{4pt}
\titlespacing*{\subsection}{0pt}{7pt}{2pt}
\titlespacing*{\subsubsection}{0pt}{5pt}{1pt}

% --- Header and Footer ---
\pagestyle{fancy}
\fancyhf{}
\renewcommand{\headrulewidth}{0.4pt}
\renewcommand{\footrulewidth}{0.4pt}
\rfoot{\small\color{monolight} Page \thepage}

% =============================================================
\begin{document}
% =============================================================

% ---- TITLE BLOCK ----
\begin{center}
  {\Huge\bfseries\color{monodark} Project X --- Lip Reading}\\[5pt]
  {\large\itshape\color{monomedium}
    Visual Speech Recognition via 3D CNN \& Bidirectional LSTM on MIRACL-VC1
  }\\[9pt]
  \noindent\rule{\linewidth}{0.4pt}\\[6pt]
\end{center}

% ---- ABSTRACT ----
\noindent\textbf{Abstract.}\quad
Project X is a deep learning based visual speech recognition system that
converts silent lip movements into text. The system applies a three-stage
preprocessing pipeline --- face detection via Haar Cascade classifiers, lip
region extraction via dlib's 68-point facial landmark model, and temporal
frame sampling with zero-padding --- to transform raw webcam or video input
into fixed-length grayscale lip sequences. A hybrid \textbf{3D CNN + Bidirectional
LSTM} model then classifies those sequences into one of 10 words or 10
phrases from the \textbf{MIRACL-VC1} dataset. This document provides a full
technical specification: dataset structure, preprocessing algorithms, model
architecture mathematics, training strategies, hyperparameter tuning,
evaluation metrics, live detection pipeline, and future roadmap.

\vspace{6pt}
\noindent\rule{\linewidth}{0.4pt}
\vspace{4pt}

% ---- SWITCH TO TWO-COLUMN BODY ----
\raggedcolumns
\begin{multicols}{2}

% ==========================================
\section{Introduction \& Motivation}
% ==========================================

Hearing-impaired individuals, noisy-environment communication, surveillance
systems, and silent speech interfaces all demand robust visual speech
recognition. Traditional audio-based ASR fails completely when audio is
absent or corrupted. Lip reading --- also called \textit{visual speech
recognition} (VSR) --- addresses this by interpreting the spatio-temporal
dynamics of lip movements.

\subsection{Challenges in Lip Reading}
\begin{itemize}
  \item \textbf{Homophones:} Words like ``p'' and ``b'' are visually identical
        (same lip shape, differ only in voicing).
  \item \textbf{Speaker variability:} Different face sizes, skin tones, and
        speaking styles.
  \item \textbf{Temporal misalignment:} Variable speaking speeds across subjects.
  \item \textbf{Occlusion \& lighting:} Real-world conditions degrade image quality.
\end{itemize}

\subsection{Project Goals}
\begin{enumerate}
  \item Build an end-to-end pipeline from webcam frames to word/phrase prediction.
  \item Achieve robust preprocessing using classical CV (Haar Cascade + dlib).
  \item Train a hybrid 3D CNN + BiLSTM model capturing spatial and temporal features.
  \item Support real-time live detection at interactive frame rates.
  \item Evaluate systematically using accuracy, precision, recall, and confusion matrices.
\end{enumerate}

% ==========================================
\section{Complete Technology Stack}
% ==========================================

\begin{center}
\small\color{monodark}
\begin{tabular}{@{}p{1.8cm} p{4.4cm}@{}}
\toprule
\textbf{Layer} & \textbf{Technology} \\
\midrule
Language      & Python 3.10+ \\
CV Library    & OpenCV (\texttt{cv2}) 4.x \\
Face Model    & Haar Cascade (OpenCV XML) \\
Landmark Det. & dlib 68-point predictor \\
Resize Util   & imutils \\
Image Proc.   & scikit-image (\texttt{skimage}) \\
DL Framework  & TensorFlow 2.x / Keras \\
Model Arch.   & 3D CNN + BiLSTM (Keras) \\
Numerics      & NumPy, Matplotlib \\
Serialization & Keras \texttt{.h5} model format \\
Dataset       & MIRACL-VC1 (Kaggle) \\
Notebooks     & Jupyter (experimentation) \\
Tools         & Git, Google Colab \\
\bottomrule
\end{tabular}
\end{center}

% ==========================================
\section{Repository File Structure}
% ==========================================

The repository is restructured into clean, purpose-separated directories
with \texttt{snake\_case} names for script compatibility.

\subsection{Layout Mapping}

\begin{center}
\scriptsize
\begin{tabular}{@{}p{3.8cm} p{3.8cm}@{}}
\toprule
\textbf{Path} & \textbf{Description} \\
\midrule
\texttt{dataset\_preprocessing/} & Preprocessing pipeline \\
\quad\texttt{01\_Face\_Extraction.ipynb} & Haar Cascade face crop \\
\quad\texttt{02\_Lip\_Extraction.ipynb} & dlib landmark lip ROI \\
\quad\texttt{03\_Train\_Test\_Split.ipynb} & Dataset splitting \\
\quad\texttt{assets/} & Model assets (\texttt{.xml}, \texttt{.dat}) \\
\texttt{model\_architecture/} & Model design \\
\quad\texttt{architectures/} & Architecture experiments \\
\quad\texttt{training\_techniques/} & Regularization \& optimizers \\
\quad\texttt{Saved Model/} & Trained \texttt{.h5} weights \\
\texttt{model\_evaluation/} & Inference \& metrics \\
\quad\texttt{Accuracy.ipynb} & Accuracy curves \\
\quad\texttt{Precision.ipynb} & Precision analysis \\
\quad\texttt{Recall.ipynb} & Recall analysis \\
\quad\texttt{Online\_testing.ipynb} & YouTube video testing \\
\quad\texttt{live\_detection\_words.py} & Words inference \\
\quad\texttt{live\_detection\_phrases.py} & Phrases inference \\
\texttt{hyperparameter\_tuning/} & Search experiments \\
\quad\texttt{grid\_search.ipynb} & Grid Search \\
\quad\texttt{random\_search.ipynb} & Random Search \\
\texttt{lip\_reading.tex} & This document \\
\texttt{README.md} & Project overview \\
\bottomrule
\end{tabular}
\end{center}

% ==========================================
\section{Dataset: MIRACL-VC1}
% ==========================================

\textbf{MIRACL-VC1} (Multi-modal Isolated Recognition of Audio and Lip
moving in Color and depth, version 1) is a benchmark dataset for visual
speech recognition.

\subsection{Dataset Composition}
\begin{itemize}
  \item \textbf{Subjects:} 15 speakers (varied demographics, accents).
  \item \textbf{Modalities:} Color video + depth video (RGB-D, from Kinect).
  \item \textbf{Utterances:} Each subject records each word/phrase 10 times.
  \item \textbf{Resolution:} $640 \times 480$ at 15 fps.
  \item \textbf{Format:} Sequences of \texttt{.png} frames per utterance.
\end{itemize}

\subsection{Vocabulary}

\subsubsection{10 Words}
\begin{codebox}
\footnotesize
\begin{verbatim}
NULL, Begin, Choose, Connection,
Navigation, Next, Previous, Start,
Stop, Hello, Web
\end{verbatim}
\end{codebox}

\subsubsection{10 Phrases}
\begin{codebox}
\footnotesize
\begin{verbatim}
Stop Navigation, Excuse me,
I am sorry, Thank you, Good bye,
I love this game, Nice to meet you,
You are welcome, How are you,
Have a good time
\end{verbatim}
\end{codebox}

\subsection{Dataset Splits}
Subjects 1--9 $\to$ \textbf{Train};
Subjects 10--15 $\to$ \textbf{Test}.
This speaker-independent split ensures the model generalises to
\textit{unseen} individuals, not just unseen recordings of known speakers.

% ==========================================
\section{Preprocessing Pipeline}
% ==========================================

The pipeline converts raw video frames into fixed-size normalised
grayscale lip sequences ready for model input.

\subsection{Stage 1: Face Detection \& Extraction}
\textbf{File:} \texttt{dataset\_preprocessing/01\_Face\_Extraction.ipynb}

\subsubsection{Tool: Haar Cascade Classifier}
OpenCV's pre-trained \texttt{haarcascade\_frontalface\_default.xml}
implements the Viola-Jones algorithm.

\subsubsection{Algorithm}
\begin{enumerate}
  \item Convert BGR frame to RGB.
  \item Run cascade with \texttt{scaleFactor=1.3}, \texttt{minNeighbors=5}.
  \item If faces detected, crop the \textbf{first} bounding box:
\end{enumerate}
\begin{codebox}
\footnotesize
\begin{verbatim}
(x, y, w, h) = faces[0]
face = image[y:y+h, x:x+w]
\end{verbatim}
\end{codebox}

\subsubsection{Viola-Jones Feature Types}
Four Haar feature types are evaluated at every scale and position
using an integral image ($O(1)$ lookup):
\begin{itemize}
  \item Two-rectangle (horizontal / vertical)
  \item Three-rectangle (centre-weighted)
  \item Four-rectangle (diagonal)
\end{itemize}

\subsubsection{Integral Image}
\begin{equation}
I_\Sigma(x,y) = \sum_{x' \leq x,\; y' \leq y} I(x', y')
\end{equation}
Rectangle sum over any region $(x_1,y_1)\to(x_2,y_2)$:
\begin{equation}
\Sigma = I_\Sigma(x_2,y_2) - I_\Sigma(x_1,y_2)
       - I_\Sigma(x_2,y_1) + I_\Sigma(x_1,y_1)
\end{equation}

\subsection{Stage 2: Lip Region Extraction}
\textbf{File:} \texttt{dataset\_preprocessing/02\_Lip\_Extraction.ipynb}

\subsubsection{Tool: dlib 68-Point Landmark Model}
\texttt{shape\_predictor\_68\_face\_landmarks.dat} is a pre-trained
regression-tree model that predicts 68 $(x, y)$ coordinates on the face.
Points 48--67 define the mouth region.

\subsubsection{Algorithm}
\begin{enumerate}
  \item Resize face crop to \texttt{width=56} px (speed optimisation).
  \item Convert to grayscale.
  \item Create a \texttt{dlib.rectangle} covering the full image.
  \item Predict 68 landmarks; convert to NumPy array via
        \texttt{face\_utils.shape\_to\_np()}.
  \item Extract mouth landmark indices using
        \texttt{FACIAL\_LANDMARKS\_IDXS['mouth']} $\to (i, j)$.
  \item Compute tight bounding rect around mouth points:
\end{enumerate}
\begin{codebox}
\footnotesize
\begin{verbatim}
(x,y,w,h) = cv2.boundingRect(
    np.array([landmarks[i:j]]))
lip = image[y-2:y+h+2, x-2:x+w+2]
\end{verbatim}
\end{codebox}
\begin{enumerate}
  \setcounter{enumi}{6}
  \item Resize lip crop to \texttt{width=500} px with \texttt{INTER\_CUBIC}.
  \item Convert to grayscale.
\end{enumerate}

\subsubsection{Landmark Index Map (mouth)}
\begin{equation}
\text{Outer lip: } 48 \to 59, \quad \text{Inner lip: } 60 \to 67
\end{equation}

\subsection{Stage 3: Frame Sampling, Resize \& Normalisation}

\subsubsection{Temporal Sampling}
Raw recordings vary in duration. A uniform sub-sampling strategy
extracts a fixed number of frames:
\begin{equation}
N_{\text{frames}} = \min\!\bigl(\text{max\_f},\; \max(\text{min\_f},\;
  \lfloor T / 10 \rfloor)\bigr)
\end{equation}
where $T$ is total frames, $\text{min\_f} = 8$, $\text{max\_f} = 28$.

\begin{equation}
\text{interval} = \max\!\left(1,\; \left\lfloor \frac{T}{N_{\text{frames}}} \right\rfloor\right)
\end{equation}

\subsubsection{Spatial Resize}
Each lip frame is resized to $100 \times 100$ px using \texttt{skimage}:
\begin{codebox}
\footnotesize
\begin{verbatim}
frame = skimage.transform.resize(lip,(100,100))
frame = (255 * frame).astype(np.uint8)
\end{verbatim}
\end{codebox}

\subsubsection{Zero Padding}
To reach a fixed sequence length of 28:
\begin{equation}
\text{pad} = 28 - N_{\text{valid}}, \quad
\mathbf{0}_{100 \times 100} \times \text{pad}
\end{equation}

\subsubsection{Normalisation}
Per-frame min-max normalisation prevents scale dominance:
\begin{equation}
\hat{x}_t = \frac{x_t - \min(x_t)}{\max(x_t) - \min(x_t)}
\end{equation}
\texttt{np.nan\_to\_num()} handles zero-division from blank padding frames.

\subsubsection{Final Input Tensor}
\begin{equation}
\mathbf{X} \in \mathbb{R}^{1 \times 28 \times 100 \times 100 \times 1}
\end{equation}
Shape: batch $\times$ frames $\times$ height $\times$ width $\times$ channels.

\columnbreak

% ==========================================
\section{Model Architecture}
% ==========================================

\textbf{Best model:} \texttt{architectures/3D\_CNN\_Bi-LSTM.ipynb} /
\texttt{Saved Model/3D\_CNN\_Bi-LSTM.h5} ($\approx$18.4 MB)

\subsection{Why 3D CNN?}
A standard 2D CNN processes each frame independently, discarding temporal
correlation. A 3D convolution applies a $k \times k \times k$ kernel across
\textit{space} (H, W) and \textit{time} (T) simultaneously:

\begin{equation}
y[t,h,w] = \sum_{\tau}\sum_{i}\sum_{j}
  K[\tau,i,j] \cdot x[t{+}\tau,\, h{+}i,\, w{+}j]
\end{equation}

\subsection{Why Bidirectional LSTM?}
Standard LSTM only sees past context. A BiLSTM processes the sequence
in both directions and concatenates hidden states:
\begin{equation}
\mathbf{h}_t = \left[\overrightarrow{\mathbf{h}}_t \;\|\; \overleftarrow{\mathbf{h}}_t\right]
\end{equation}
Critical for lip reading because some phoneme boundaries are only
disambiguated by \textit{future} lip context (co-articulation).

\subsection{Layer-by-Layer Architecture}

\begin{center}
\scriptsize
\begin{tabular}{@{}p{2.4cm} p{1.8cm} p{1.5cm}@{}}
\toprule
\textbf{Layer} & \textbf{Output Shape} & \textbf{Params} \\
\midrule
Input              & (28,100,100,1)     & 0 \\
Conv3D(32,3,ReLU)  & (26,98,98,32)      & 896 \\
MaxPool3D(2,2)     & (13,49,49,32)      & 0 \\
Conv3D(64,3,ReLU)  & (11,47,47,64)      & 55,360 \\
MaxPool3D(2,2)     & (5,23,23,64)       & 0 \\
Conv3D(128,3,ReLU) & (3,21,21,128)      & 221,312 \\
MaxPool3D(2,2)     & (1,10,10,128)      & 0 \\
Reshape            & (1,12800)           & 0 \\
BiLSTM(256)        & (1,512)             & $\sim$27.4M \\
Flatten            & (512,)              & 0 \\
Dropout(0.5)       & (512,)              & 0 \\
Dense(256, ReLU)   & (256,)              & 131,328 \\
Dense(11, Softmax) & (11,)               & 2,827 \\
\midrule
\multicolumn{2}{l}{\textbf{Total Trainable}} & $\approx$27.8M \\
\bottomrule
\end{tabular}
\end{center}

\subsection{3D CNN Block Equations}

For layer $l$ with $F$ filters, kernel $\mathbf{K}^l$:
\begin{equation}
\mathbf{A}^l = \text{ReLU}\!\left(
  \mathbf{K}^l * \mathbf{A}^{l-1} + \mathbf{b}^l
\right), \quad \text{ReLU}(z) = \max(0,z)
\end{equation}

MaxPool3D halves spatial \& temporal dimensions with no parameters.

\subsection{LSTM Cell Equations}

\begin{align}
\mathbf{f}_t &= \sigma(\mathbf{W}_f \mathbf{h}_{t-1} + \mathbf{U}_f \mathbf{x}_t + \mathbf{b}_f) \\
\mathbf{i}_t &= \sigma(\mathbf{W}_i \mathbf{h}_{t-1} + \mathbf{U}_i \mathbf{x}_t + \mathbf{b}_i) \\
\mathbf{o}_t &= \sigma(\mathbf{W}_o \mathbf{h}_{t-1} + \mathbf{U}_o \mathbf{x}_t + \mathbf{b}_o) \\
\tilde{\mathbf{c}}_t &= \tanh(\mathbf{W}_c \mathbf{h}_{t-1} + \mathbf{U}_c \mathbf{x}_t + \mathbf{b}_c) \\
\mathbf{c}_t &= \mathbf{f}_t \odot \mathbf{c}_{t-1} + \mathbf{i}_t \odot \tilde{\mathbf{c}}_t \\
\mathbf{h}_t &= \mathbf{o}_t \odot \tanh(\mathbf{c}_t)
\end{align}

$\mathbf{f}_t$ (forget), $\mathbf{i}_t$ (input), $\mathbf{o}_t$ (output)
are the three gates; $\sigma$ is sigmoid; $\odot$ is elementwise product.

\subsection{Output: Softmax}

\begin{equation}
P(\hat{y} = k \mid \mathbf{x}) = \frac{e^{z_k}}{\displaystyle\sum_{j=1}^{K} e^{z_j}},
\quad K \in \{11, 10\}
\end{equation}
Prediction: $\hat{y} = \arg\max_k P(k \mid \mathbf{x})$.

% ==========================================
\section{Training \& Regularization}
% ==========================================

\subsection{Loss: Categorical Cross-Entropy}

\begin{equation}
\mathcal{L} = -\sum_{k=1}^{K} y_k \log \hat{p}_k
\end{equation}

$y_k \in \{0,1\}$ is one-hot label; $\hat{p}_k$ is predicted probability.

\subsection{Optimizers}

\subsubsection{Adam}
Adaptive Moment Estimation maintains first and second moment estimates:
\begin{align}
m_t &= \beta_1 m_{t-1} + (1-\beta_1)g_t \\
v_t &= \beta_2 v_{t-1} + (1-\beta_2)g_t^2 \\
\theta_t &= \theta_{t-1} - \frac{\eta}{\sqrt{\hat{v}_t}+\epsilon}\hat{m}_t
\end{align}
Default: $\eta=10^{-3}$, $\beta_1=0.9$, $\beta_2=0.999$, $\epsilon=10^{-8}$.

\subsubsection{RMSprop}
\begin{equation}
E[g^2]_t = \rho E[g^2]_{t-1} + (1{-}\rho)g_t^2, \quad
\theta_t = \theta_{t-1} - \frac{\eta}{\sqrt{E[g^2]_t + \epsilon}} g_t
\end{equation}
$\rho=0.9$ typical. Adam outperformed RMSprop on this task.

\subsection{Dropout}
Each neuron zeroed with probability $p$ during training:
\begin{equation}
\tilde{\mathbf{h}} = \mathbf{m} \odot \mathbf{h}, \quad
m_i \sim \text{Bernoulli}(1-p)
\end{equation}
Rate $p = 0.5$ used in Dense block (inverted dropout at inference).

\subsection{L1 \& L2 Regularisation}

\subsubsection{L1 (Lasso) --- encourages sparsity}
\begin{equation}
\mathcal{L}_{\text{total}} = \mathcal{L}_{\text{CE}} + \lambda \sum_i |w_i|
\end{equation}

\subsubsection{L2 (Ridge / Weight Decay)}
\begin{equation}
\mathcal{L}_{\text{total}} = \mathcal{L}_{\text{CE}} + \frac{\lambda}{2} \sum_i w_i^2
\end{equation}
Gradient contribution $\lambda w_i$ uniformly shrinks all weights.
L2 is smoother and more stable for deep networks.

\subsection{Early Stopping}
Monitors \texttt{val\_loss}; halts when no improvement for \texttt{patience} epochs:
\begin{equation}
\text{val\_loss}_{e} \geq \text{val\_loss}_{e - p}
\quad \forall\; p \in [1, \text{patience}]
\end{equation}
Best weights from $e^*$ are restored. \texttt{patience=10} used.

\subsection{Data Augmentation}
Applied to lip frame sequences to increase effective dataset size:
\begin{itemize}
  \item Horizontal flip: $x \leftarrow W - 1 - x$
  \item Brightness jitter: pixel scaling $\in [0.8, 1.2]$
  \item Gaussian noise: $x \leftarrow x + \mathcal{N}(0,\sigma^2)$
  \item Random crop: slight spatial shifts within lip region
\end{itemize}

% ==========================================
\section{Hyperparameter Tuning}
% ==========================================

\subsection{Search Space}

\begin{center}
\scriptsize
\begin{tabular}{@{}p{1.9cm} p{1.3cm} p{2.3cm}@{}}
\toprule
\textbf{Parameter} & \textbf{Type} & \textbf{Range / Options} \\
\midrule
Learning rate   & Continuous & $[10^{-4}, 10^{-2}]$ \\
Dropout rate    & Continuous & $[0.3, 0.6]$ \\
LSTM units      & Discrete   & $\{128, 256, 512\}$ \\
Dense units     & Discrete   & $\{128, 256\}$ \\
L2 $\lambda$    & Continuous & $[10^{-5}, 10^{-2}]$ \\
Batch size      & Discrete   & $\{16, 32\}$ \\
\bottomrule
\end{tabular}
\end{center}

\subsection{Grid Search}
Exhaustively evaluates every combination of a discrete grid.
For $n$ parameters with $k$ values each: $O(k^n)$ evaluations.

\subsection{Random Search}
Samples $N$ random combinations. Empirically finds good configs
with far fewer evaluations than grid search (Bergstra \& Bengio, 2012):
\begin{equation}
P(\text{finding } \epsilon\text{-optimal in }N\text{ trials})
= 1 - (1 - \epsilon)^N
\end{equation}

\subsubsection{Best Configuration Found}
\begin{codebox}
\footnotesize
\begin{verbatim}
lr          = 1e-3    (Adam)
dropout     = 0.5
lstm_units  = 256     (BiLSTM -> 512)
dense_units = 256
l2_lambda   = 1e-4
batch_size  = 32
\end{verbatim}
\end{codebox}

% ==========================================
\section{Model Evaluation}
% ==========================================

\subsection{Accuracy}
\begin{equation}
\text{Accuracy} = \frac{TP + TN}{TP + TN + FP + FN}
\end{equation}
Training/validation accuracy curves per epoch reveal overfitting when
they diverge significantly.

\subsection{Precision}
\begin{equation}
\text{Precision}_k = \frac{TP_k}{TP_k + FP_k}
\end{equation}
High precision: model predicts class $k$ only when confident.

\subsection{Recall (Sensitivity)}
\begin{equation}
\text{Recall}_k = \frac{TP_k}{TP_k + FN_k}
\end{equation}
High recall: few actual class-$k$ samples are missed.

\subsection{F1 Score}
\begin{equation}
F_1 = \frac{2 \cdot \text{Precision} \cdot \text{Recall}}
           {\text{Precision} + \text{Recall}}
\end{equation}

\subsection{Confusion Matrix}
A $K \times K$ matrix $\mathbf{C}$ where $C_{ij}$ counts samples of
true class $i$ predicted as class $j$.
Diagonal = correct; off-diagonal = errors.
Commonly confused pairs reveal visually ambiguous (homophenic) words.

\subsection{Online Testing}
\textbf{File:} \texttt{model\_evaluation/Online\_testing.ipynb}

Model evaluated on video clips from YouTube (real-world unseen speakers)
to test generalisation beyond the MIRACL-VC1 distribution.

% ==========================================
\section{Live Detection System}
% ==========================================

\textbf{Files:} \texttt{live\_detection\_words.py},
\texttt{live\_detection\_phrases.py}

\subsection{End-to-End Pipeline}

\begin{enumerate}
  \item \textbf{Capture:} Open webcam via \texttt{cv2.VideoCapture(0)}.
  \item \textbf{Control:} \texttt{SPACE} toggles recording; \texttt{ESC}
        exits. Frames buffered into list.
  \item \textbf{Frame sampling:} Apply Eq.\,(8) using $T = |\text{frames}|$.
  \item \textbf{Face extraction:} Haar Cascade per sampled frame.
  \item \textbf{Lip extraction:} dlib 68-landmark ROI crop.
  \item \textbf{Resize \& pad:} $100 \times 100$ + zero-pad to 28 frames.
  \item \textbf{Normalise:} Min-max per frame (Eq.\,10).
  \item \textbf{Reshape:} $(1, 28, 100, 100, 1)$.
  \item \textbf{Inference:} \texttt{loaded\_model.predict(sequence)}.
  \item \textbf{Output:} Argmax class + confidence percentages printed.
\end{enumerate}

\subsection{Key Design Decisions}
\begin{itemize}
  \item \textbf{Flip display:} \texttt{cv2.flip(frame,1)} mirrors feed
        (selfie-camera feel).
  \item \textbf{Buffer copy:} \texttt{frames.append(frame.copy())} avoids
        reference aliasing from rolling capture buffer.
  \item \textbf{Skip on failure:} Frames where no face or lips are detected
        are skipped gracefully, making pipeline robust to occlusions.
  \item \textbf{Error handling:} All extraction steps wrapped in try/except
        with informative console output.
\end{itemize}

% ==========================================
\section{Architecture Variants Explored}
% ==========================================

\begin{center}
\scriptsize
\begin{tabular}{@{}p{2.5cm} p{2.2cm} p{1.0cm}@{}}
\toprule
\textbf{Architecture} & \textbf{Description} & \textbf{Val Acc} \\
\midrule
3D CNN              & 3D Conv, no RNN       & Low \\
3D CNN (Scratch)    & Custom, no pretrain   & Low \\
3D CNN + LSTM       & Unidirectional LSTM   & Medium \\
3D CNN + BiLSTM     & Bidirectional LSTM    & \textbf{Best} \\
\bottomrule
\end{tabular}
\end{center}

The BiLSTM consistently outperformed unidirectional LSTM due to its
ability to leverage future lip context (co-articulation effects),
which is crucial for distinguishing visually similar words.

% ==========================================
\section{Conclusion \& Future Scope}
% ==========================================

Project X demonstrates that combining classical CV preprocessing
(Haar Cascade + dlib) with a deep hybrid architecture (3D CNN + BiLSTM)
produces a functional real-time lip reading system. The modular, refactored
repository structure facilitates further research and extension.

\subsection{Future Directions}
\begin{itemize}
  \item \textbf{Sentence-level reading:} Use CTC or attention decoders
        for continuous sentence recognition beyond isolated words.
  \item \textbf{Multilingual support:} Train on multilingual corpora
        (LRS2, LRS3) for broader applicability.
  \item \textbf{Transformer backbone:} Replace BiLSTM with a Temporal
        Transformer (e.g., LipFormer) for better long-range dependency.
  \item \textbf{Larger datasets:} Train on LRW (500 words) or LRS3
        (100k+ sentences) for production-grade generalisation.
  \item \textbf{Audio-Visual fusion:} Combine lip features with audio
        for robust AV-ASR in noisy environments.
  \item \textbf{Edge deployment:} Quantize and export to ONNX / TFLite
        for real-time inference on embedded devices.
  \item \textbf{GUI interface:} Build live subtitle overlay for
        accessibility applications.
\end{itemize}

\end{multicols}
\end{document}
